Edmond Halley (1656-1742) va estudiar els registres històrics d'observacions de cometes dels segles \textsc{xvi} i \textsc{xvii}, i va concloure que diverses d'aquestes observacions corresponien al mateix objecte. Es tractava d'un cometa periòdic gran i brillant que orbitava al voltant del Sol i que havia estat observat per darrera vegada l'any 1682. Actualment aquest cometa és conegut com el cometa Halley.

\begin{apartat}{1,25}
Feu un dibuix esquemàtic de l'òrbita el·líptica del cometa Halley, indicant-ne el periheli i l'afeli. Calculeu el semieix major de l'el·lipse i el període orbital del cometa. Raoneu si Edmond Halley va poder observar el cometa que porta el seu nom més d'una vegada.
\end{apartat}

\begin{apartat}{1,25}
Raoneu en quin punt de l'òrbita el cometa assoleix la velocitat màxima. A partir de les dades que es faciliten, calculeu la velocitat del cometa Halley quan passa pel periheli de l'òrbita, indicant clarament els principis físics que heu emprat.
\end{apartat}

\dades{%
  $G = 6,67\times 10^{-11}\ \mathrm{N\,m^2\,kg^{-2}}$.\\
  Massa del Sol: $M_\mathrm{S} = 1,99\times 10^{30}\ \mathrm{kg}$.\\
  Distància del periheli al centre del Sol: $r_\mathrm{P} = 0,575\ \mathrm{ua}$.\\
  Velocitat a l'afeli: $v_\mathrm{A} = 896\ \mathrm{m\,s^{-1}}$.\\
  Distància de l'afeli al centre del Sol: $r_\mathrm{A} = 35,3\ \mathrm{ua}$.\\
  $1\ \mathrm{ua} = 1,50\times 10^{11}\ \mathrm{m}$.}
